% Performance: large document with many sections, code blocks, and diagrams.
% Compile time with two passes should stay under the configured limit.
% !TEX program = xelatex
\documentclass[code=listings]{SUSTechHomework}

\title{Large Document Performance Test}
\coursecode{PERF001}
\coursename{Performance Testing}
\author{Perf Tester}
\sid{99990001}
\semester{Spring 2026}

\begin{document}
\maketitle

% --- 20 sections with mixed content to simulate a real large report ---

\section{Introduction}
Performance test document. Lorem ipsum dolor sit amet, consectetur adipiscing elit.

\section{Background}
Sed do eiusmod tempor incididunt ut labore et dolore magna aliqua.

\begin{sustechnote}
Background note block. Ut enim ad minim veniam.
\end{sustechnote}

\section{Related Work}
Quis nostrud exercitation ullamco laboris nisi ut aliquip ex ea commodo consequat.

\begin{sustechcode}{Python}
# Section 3 code block
def related_work():
    return "processing"
\end{sustechcode}

\section{Methodology}
Duis aute irure dolor in reprehenderit in voluptate velit esse cillum dolore.

\begin{table}[htbp]
  \centering
  \begin{tabular}{lll}
    \toprule
    \sustechheaderrow
    \sustechtableheader{Method} & \sustechtableheader{Precision} & \sustechtableheader{Recall} \\
    \midrule
    Baseline & 0.82 & 0.78 \\
    Ours     & 0.91 & 0.88 \\
    \bottomrule
  \end{tabular}
  \caption{Methodology comparison table}
\end{table}

\section{Implementation}
Excepteur sint occaecat cupidatat non proident.

\begin{sustechcode}{C++}
// Implementation code
#include <vector>
std::vector<int> data = {1, 2, 3, 4, 5};
\end{sustechcode}

\section{Experiments}
Sunt in culpa qui officia deserunt mollit anim id est laborum.

\begin{sustechfiguregroup}
  \begin{sustechsubfigure}[0.48\linewidth][Experiment A]
    \sustechplaceholdergraphic[0.92\linewidth][3cm]
  \end{sustechsubfigure}
  \sustechfloatsep
  \begin{sustechsubfigure}[0.48\linewidth][Experiment B]
    \sustechplaceholdergraphic[0.92\linewidth][3cm]
  \end{sustechsubfigure}
  \caption{Experimental results comparison}
\end{sustechfiguregroup}

\section{Results}
Lorem ipsum dolor sit amet.

\section{Analysis}
Consectetur adipiscing elit.

\begin{sustechwarning}[title=Important Result]
Analysis warning block — key finding.
\end{sustechwarning}

\section{Discussion}
Sed do eiusmod tempor incididunt.

\section{Ablation Study}
Ut labore et dolore magna aliqua.

\begin{sustechcode}{json}
{
  "ablation": true,
  "component": "encoder",
  "delta": -0.03
}
\end{sustechcode}

\section{Limitations}
Quis nostrud exercitation ullamco laboris.

\section{Future Work}
Nisi ut aliquip ex ea commodo consequat.

\section{Related Applications}
Duis aute irure dolor in reprehenderit.

\begin{figure}[htbp]
  \centering
  \begin{sustechpipelinediagram}
    \sustechphasebox{s1}{at={(0,0)}}{Stage 1}
    \sustechphasebox{s2}{right=of s1}{Stage 2}
    \sustechphasebox{s3}{right=of s2}{Stage 3}
  \end{sustechpipelinediagram}
  \caption{Future work pipeline}
\end{figure}

\section{Conclusion}
In voluptate velit esse cillum dolore eu fugiat nulla pariatur.

\section{Acknowledgements}
Excepteur sint occaecat cupidatat non proident sunt in culpa.

\section{References}
Qui officia deserunt mollit anim id est laborum.

\section{Appendix A}
Supplementary material.

\begin{sustechcode}{yaml}
appendix:
  section: A
  type: supplementary
\end{sustechcode}

\section{Appendix B}
Additional experiments.

\begin{itemize}
  \item Additional result one
  \item Additional result two
  \item Additional result three
\end{itemize}

\begin{enumerate}
  \item Step one of appendix procedure
  \item Step two of appendix procedure
\end{enumerate}

\end{document}
